\documentclass[a4paper,11pt]{article}
\usepackage[utf8]{inputenc}
\usepackage[german]{babel}
\usepackage{geometry}
\usepackage{graphicx}
\usepackage{parskip}
\usepackage{titlesec}
\usepackage{hyperref}

\setcounter{secnumdepth}{4}

\titleformat{\paragraph}
{\normalfont\normalsize\bfseries}{\theparagraph}{1em}{}
\titlespacing*{\paragraph}
{0pt}{3.25ex plus 1ex minus .2ex}{1.5ex plus .2ex}

\geometry{a4paper, margin=2.5cm}
\graphicspath{{images/}}

\begin{document}

\begin{table}[h]
	\centering
	\footnotesize
	\begin{tabular}{| p{1.5cm} | p{2.5cm} | p{2cm} | p{2cm} | p{2cm} |}
		\hline
		                            & Verantwortliche Teammitglieder (Organisation, Anlegen der Dokumentenversion, Moderation, Abgabe)         & Anwesende während der Meetings                                                                           & (Online-) Beiträge zum Inhalt durch                                                                      & Korrektur-gelesen durch                                                                                  \\
		\hline
		Dokument Sprint 3 Version 1 & Paul Schröder, Eric Buthmann, Yanik Buthmann, Mathis Behrendt, Adrian Gode, Niklas Krüger, Kevin Zellmer & Paul Schröder, Eric Buthmann, Yanik Buthmann, Mathis Behrendt, Adrian Gode, Niklas Krüger, Kevin Zellmer & Paul Schröder, Eric Buthmann, Yanik Buthmann, Mathis Behrendt, Adrian Gode, Niklas Krüger, Kevin Zellmer & Paul Schröder, Eric Buthmann, Yanik Buthmann, Mathis Behrendt, Adrian Gode, Niklas Krüger, Kevin Zellmer \\
		\hline
	\end{tabular}
\end{table}

\section{Backlogs und Planung}

\subsection{Produktbacklog vor und nach dem Sprint}


\subsection{Sprint Planung}


\section{Verfeinerte Architektur}


\section{Implementation und Tests}
\subsection{Updates zu den genutzen Technologien}
An den genutzten Technologien wurden keine Änderungen vorgenommen.


\subsection{Dokumentation der Codequalität}


\subsection{Tracing}


\subsection{Laufender Prototyp}


\subsubsection*{Installations- und Kompilieranweisung}

Im Folgenden wird beschrieben, wie man den Prototypen kompilieren kann:

\textbf{Voraussetzungen:}
\begin{itemize}
	\item Java 21 oder höher (getestet unter Windows 11 und Ubuntu 24.04)
	\item Apache Maven 3.9 oder höher
	\item SQLite (keine separate Installation nötig -- der JDBC-Treiber ist im Build enthalten)
\end{itemize}

\textbf{Schritt 1 -- Repository klonen:}
\begin{verbatim}
git clone https://git.informatik.uni-rostock.de/softwaretechnik-ws-2025-26/
    uebung3/team3/projekt-team3-uebung3.git
cd projekt-team3-uebung3
\end{verbatim}

\textbf{Schritt 2 -- Konfigurationsdatei anpassen:}

Die Datei \texttt{src/main/resources/application.properties} enthält den Parameter \texttt{database.url}.
Dieser muss auf einen gültigen Pfad zur SQLite-Datenbankdatei gesetzt werden, z.\,B.:
\begin{verbatim}
database.url=jdbc:sqlite:/home/user/scrumboard.db
\end{verbatim}
Die Datenbankdatei wird beim ersten Start automatisch angelegt.

\textbf{Schritt 3 -- Projekt bauen:}
\begin{verbatim}
mvn package
\end{verbatim}

\textbf{Schritt 4 -- Anwendung starten:}
\begin{verbatim}
java -jar target/scrumboard-1.0-SNAPSHOT.jar
\end{verbatim}

Die Anwendung ist anschließend im Browser unter localhost erreichbar.

\subsection{Abweichung Sprintplanung}


\end{document}

